% Generated draft: requires mathematical and source review.
\documentclass[11pt]{article}
\usepackage[margin=25mm]{geometry}
\usepackage{fontspec}
\setmainfont{Latin Modern Roman}
\usepackage{amsmath,amssymb,mathtools,unicode-math}
\setmathfont{Latin Modern Math}
\usepackage{enumitem,needspace,xurl,hyperref}
\setlength{\parindent}{0pt}
\setlength{\parskip}{5pt}
\setlength{\emergencystretch}{3em}
\newcommand{\N}{\mathbb{N}}
\newcommand{\Z}{\mathbb{Z}}
\newcommand{\Q}{\mathbb{Q}}
\newcommand{\R}{\mathbb{R}}
\newcommand{\C}{\mathbb{C}}
\newcommand{\F}{\mathbb{F}}
\newcommand{\A}{\mathbb{A}}
\newcommand{\PP}{\mathbb{P}}
\newcommand{\E}{\mathbb{E}}
\newcommand{\eps}{\varepsilon}
\newcommand{\dd}{\,\mathrm{d}}
\begin{document}
\section*{UnsolvedMath ID: 30005311}
\textbf{Source problem code:} catalogue entry
\subsection{Definitions and mathematical statement}
\textbf{Scope and terminology.} The supplied record concerns historical and archival research on at least one of Oswald Veblen's lecture notebooks on the foundations of geometry associated in the excerpt with 1901. A \emph{scriptor} means a person who writes the manuscript text; identifying a scriptor does not by itself identify the lecturer or the originator of the recorded ideas. The excerpt distinguishes the voice of the note taker from that of the speaker or written text being recorded. It recommends investigating the writing practices through which concepts are formed, transmitted, and appropriated, while cautioning that some general questions about lecture notes may be less relevant to Veblen's notebook.

\textbf{Recoverable historical inquiry.} For the manuscript intended by the source, investigate whether there are changes in the pens used and/or in handwriting, and determine when the manuscript was written. The surrounding discussion also asks who wrote the notes and whether more than one scriptor participated. These are research questions, not assertions that changes or multiple scriptors exist.

\textbf{Unresolved formulation.} The excerpt does not identify the manuscript precisely, specify which portions are to be compared, or establish whether the opening questions concern exactly the notebook subsequently associated with 1901. That association does not establish a precise composition date for every part of the manuscript. The text ends with an incomplete question about the lecturer, so its intended content cannot be recovered. No mathematical objects, hypotheses, or proposition suitable for a mathematical problem statement are supplied. The record's supplied statement-verification note classifies the passage as a historical and archival inquiry rather than an open mathematical problem, and its supplied statement status is \texttt{unrecoverable}. Accordingly, this draft preserves the recoverable inquiry without inventing a mathematical formulation.
\subsection{Short English statement}
Investigate changes in pens or handwriting, the date of composition, and the identity and possible multiplicity of writers in the manuscript discussed alongside Veblen's geometry lecture notebooks associated with 1901. The manuscript is not precisely identified and the excerpt is truncated. This is a historical inquiry with no recoverable mathematical proposition; the supplied \texttt{open} label is a curation label, not an independently verified mathematical status.
\subsection{Sources}
\begin{itemize}
\item Local UnsolvedMath registry: supplied complete record, permanent ID \(30005311\), title \emph{Handwriting and Chronology in Veblen's Notebooks}, exact source problem code \texttt{catalogue entry}, in the Oberwolfach Reports Open Problems collection. No individual proposer is supplied; the registry gives proposed year 2022.
\item Supplied report citation: \emph{History of mathematics through collaboration: Toward a composite portrait of Oswald Veblen} (2023), Oberwolfach Reports, pp.\ 3171--3214, DOI \texttt{10.4171/owr/2022/56}. Supplied source URL and literature-source link labeled ``original OWR report'': \url{https://doi.org/10.4171/owr/2022/56}. The record describes Oberwolfach Reports as documenting mathematical workshops and their open-problem discussions; that provenance does not establish that this particular excerpt states a mathematical problem.
\item Supplied literature-source link labeled ``Library of Congress primary finding aid'': \url{https://findingaids.loc.gov/repositories/19/resources/4615}.
\item Supplied literature assessment, checked \textbf{2026-08-22}: the record labels the inquiry \texttt{open} and reports that the archive confirms the existence of relevant lecture notebooks, while the handwriting-change and dating questions remain unsettled by the available archival description or scholarship. This is supplied dated curation derived from the OWR record, not an independently verified assessment. The record itself requests independent verification of time-sensitive status and cited claims.
\item The excerpt attributes methodological observations to its reference [18] and contains a footnote marker 6, but supplies neither the bibliographic identity of [18] nor the footnote text. These attributions cannot be completed from the supplied evidence. No authoritative citation establishing a mathematical problem is supplied.
\end{itemize}
\end{document}
